\documentclass{shumo-cumcm2025}
\begin{document}
\papertitle{基于某模型的某问题研究}
\begin{cnabstract}
本文针对……问题，建立……模型。针对问题一，……；针对问题二，……；针对问题三，……。给出核心定量结果、检验结论和适用边界，摘要不超过一页。
\end{cnabstract}
\keywords{关键词一；关键词二；关键词三}
\clearpage
\section{问题重述}
用原创语言说明背景、目标和约束。
\section{问题分析}
说明总体路线和各小问关键矛盾，不在此处提前堆放结果。
\section{模型假设}
列出与建模直接相关且有依据的假设。
\section{符号说明}
使用三线表列出全局符号、含义和单位。
\section{模型建立与求解}
\subsection{问题一}
先定义变量，再建立模型。仅核心公式编号：
\keyequation{eq:example}{y=f(x;\bm\theta).}
在正文中引用式~\eqref{eq:example}，并给出求解、结果分析和直接回答。
\section{模型检验与讨论}
给出一致性、误差、敏感性或对照检验，并保留不利结果。
\section{模型评价}
具体说明优点、局限和改进方向。
\section{结论}
直接汇总各问结论和适用范围。
\begin{thebibliography}{9}
\bibitem{sample} 作者. 题名[J]. 期刊, 年, 卷(期): 页码.
\end{thebibliography}
\clearpage
\appendix
\section{支撑材料文件列表}
列出论文对应的全部程序、数据和图形源文件。
\section{完整源程序}
\lstinputlisting{example.py}
\end{document}
